\documentclass[11pt]{article}
\usepackage[margin=1in]{geometry}
\usepackage{amsmath,amssymb,amsthm,mathtools,booktabs,microtype}
\usepackage[hidelinks]{hyperref}
\newtheorem{theorem}{Theorem}
\newtheorem{lemma}[theorem]{Lemma}
\newcommand{\E}{\mathbb E}
\newcommand{\TV}{\operatorname{TV}}
\newcommand{\diag}{\operatorname{diag}}
\newcommand{\Forr}{\operatorname{Forr}_4}
\title{Fourfold forrelation at $N=4096$:\\a construction and a conditional lower bound}
\author{Amir H. Safavi-Naeini}
\date{}
\begin{document}
\maketitle
\begin{abstract}
Three single photons, each crossing two phase masks, solve the
fourfold-forrelation promise with hard dose six and error $81/256$.
We reduce a proposed parallel lower bound to explicit occurrence-kernel
inequalities and a finite spectral calculation. The spectral arithmetic
is reproducible, but supplied coefficient derivations contain errors:
the finite lower bound remains conditional. We also correct its probe
restriction from physical photon number to parity-support size.
\end{abstract}

\section{The task and the status of the results}
For $N=2^d$, index coordinates by $\mathbb F_2^d$ and set
$H_{ij}=N^{-1/2}(-1)^{i\cdot j}$,
$u=N^{-1/2}(1,\ldots,1)^{\mathsf T}$.
Four binary phase masks $D_b=\diag(x^{(b)})$ define
\begin{equation}
 F(x)=\Forr(x)=u^{\mathsf T}D_1HD_2HD_3HD_4u .
 \label{eq:task}
\end{equation}
Decide $F\ge1/4$ versus $F\le-1/4$ with worst-case error at most $1/3$.
This is a variant of the query problem of
Aaronson and Ambainis~\cite{AaronsonAmbainis2015}.
Dose counts every photon-mask traversal and bounds the entire support
of the signal state, not its mean or a postselected cost.

The two-pass upper bound below is unconditional.
The proposed finite obstruction concerns one parallel interrogation:
a fixed probe passes once through the direct-sum mask and is measured
by an arbitrary POVM. References and idlers bypass the masks.
For a Fock occupation $\nu$, define
\[
 S_b(\nu)=\{i:\nu_{b,i}\text{ is odd}\},\qquad
 Q|\nu\rangle=\left(\sum_b|S_b(\nu)|\right)|\nu\rangle.
\]
The balanced coefficient ledger requires
\begin{equation}
 [\rho,Q]=0,\qquad
 \rho\text{ supported on }\widehat N_{\rm sig}\le6.
 \label{eq:number}
\end{equation}
Thus the probe is block diagonal in parity-support size. Repeated modes
and idlers are allowed. Fixed physical photon number alone is insufficient:
the superposition of three photons in one mode of block one and one
photon in each of blocks two, three and four has parity sizes one and
three. The components acquire different phase characters.

\begin{theorem}[Conditional finite obstruction]\label{thm:finite}
Assume the complete-kernel bounds \eqref{eq:coefficient} hold with the
recorded coefficients. At $N=4096$ there are laws $P_+,P_-$ on the
respective promises such that every parallel probe satisfying
\eqref{eq:number} has output laws obeying
\begin{equation}
 \TV(T_+,T_-)\le0.260969224792207925.
 \label{eq:total}
\end{equation}
Its equal-prior Bayes error is at least
$0.369515387603896037>1/3$.
\end{theorem}

The hypothesis is not established by the recorded arithmetic.
Independent review found false local-Walsh and four-cubic intermediate
bounds in the supplied derivations. Section~\ref{sec:ledger} states
exactly what remains to be proved; \path{../AUDIT.md} gives the
counterexamples. No finite parallel or adaptive separation is claimed
unconditionally by this certificate.

\section{The two-pass protocol}
Prepare $(|0\rangle+|1\rangle)u/\sqrt2$.
In arm zero apply $D_1,H,D_2$; in arm one apply $D_4,H,D_3,H$.
The resulting state is $(|0\rangle L_x+|1\rangle R_x)/\sqrt2$, where
\[
 L_x=D_2HD_1u,\qquad R_x=HD_3HD_4u.
\]
Each arm crosses exactly two masks. Recombining the paths gives
\begin{equation}
 \Pr(X=\pm1\mid x)=\frac{1\pm\langle L_x,R_x\rangle}{2}
                 =\frac{1\pm F(x)}2.
 \label{eq:flag}
\end{equation}
Discard the mode label and retain the output port.
On the promise, each flag is correct with probability at least $5/8$.
Three independent flags and majority decoding have error at most
\[
 (3/8)^3+3(5/8)(3/8)^2=81/256<1/3,
\]
and deterministic dose $3\cdot2=6$. Equality holds at either promise
boundary. No outcome is discarded. This construction works for every $N$.

\section{The hard inputs}
Set $q=64$, so $N=q^2$. Identify a coordinate with $(a,b)\in[q]^2$,
and use row-major flattening, so $H_N=H_q\otimes H_q$.
Write $S_q=\sqrt q\,H_q$.
For a uniform permutation $\pi$ and independent uniform signs $\sigma_b$,
define arrays
\begin{equation}
 L_{a,b}=S_q(a,\pi(b))\sigma_b,\qquad
 R_{\pi(b),a}=\sigma_bS_q(b,a).
 \label{eq:lr}
\end{equation}
Their useful property is simply
\begin{equation}
                         H_N L=R,\qquad H_N R=L.
 \label{eq:hlr}
\end{equation}
Indeed, at coordinate $(c,d)$,
\[
 (H_NL)_{c,d}
 =\frac1q\sum_{a,b}S_q(c,a)S_q(d,b)S_q(a,\pi(b))\sigma_b
 =S_q(d,b_c)\sigma_{b_c}=R_{c,d},
\]
where $b_c=\pi^{-1}(c)$. The second identity follows from $H_N^2=I$.

Take three independent pairs $(L_j,R_j)$ and set
\begin{equation}
 z=(L_1,\ R_1\odot L_2,\ R_2\odot L_3,\ R_3).
 \label{eq:plant}
\end{equation}
Starting on the right of \eqref{eq:task}, \eqref{eq:hlr} successively
cancels $R_3,L_3,R_2,L_2,R_1,L_1$ and leaves $u$.
Thus $F(z)=1$ exactly. Negating the first block changes the value to $-1$.

Now independently flip every coordinate with probability $3/25$.
Equivalently, multiply it by a sign $\eta$ with mean
\[
                         \beta=19/25.
\]
Let $\widetilde P_+$ be the resulting positive law and
$\widetilde P_-$ its first-block reflection.
Then $\E_{\widetilde P_+}F=\beta^4$.
Finally condition each law onto its own promise to obtain $P_+,P_-$.

\begin{lemma}[Cost of conditioning]\label{lem:promise}
The conditioning loss is at most
$0.002225128406630703$.
\end{lemma}
\begin{proof}
Fix an exact plant $z$. Reveal the four noisy blocks in order from
left to right and form the Doob martingale of $F$.
When block $b$ is exposed, the future $4-b$ masks contribute their
means, giving a factor $\beta^{4-b}$. By \eqref{eq:hlr}, the remaining
deterministic right-hand chain is a sign vector divided by $\sqrt N$.
The left-hand chain has norm one.
The sum of squared coefficients of the newly exposed signs is therefore
$\beta^{2(4-b)}/N$.

A sign in $[-1,1]$ has centered subgaussian proxy at most one by
Hoeffding's lemma. Conditional independence within each newly exposed
block gives the martingale proxy
\[
 V_\beta=\frac{1+\beta^2+\beta^4+\beta^6}{N}
        =\frac{(1+\beta^2)(1+\beta^4)}N .
\]
With $g=\beta^4-1/4$ and $x=g^2/(2V_\beta)$, Chernoff gives a
failure probability at most $e^{-x}$. This holds for every plant.
Reflection gives the same bound on the negative side.
Conditioning changes each input law in total variation by at most
$e^{-x}$. Data processing and the triangle inequality therefore add
at most $2e^{-x}$ to the output total variation.
The verifier evaluates the exact rational upper
\[
 2e^{-x}\le2(1+x/4096)^{-4096}
          <0.002225128406630703 .
\]
\end{proof}

\section{From parity supports to one matrix}\label{sec:matrix}
The mask acts on an occupation only through its parity support:
$U_x|\nu\rangle=\chi_{S(\nu)}(x)|\nu\rangle$, where
$\chi_S(x)=\prod_b\prod_{i\in S_b}x_i^{(b)}$.
By \eqref{eq:number}, decompose the input into fixed-$Q$ sectors and
purify each sector separately. Group all Fock amplitudes with support
$S$ into a vector $\psi_S$, including idlers. These vectors are mutually
orthogonal; put $p_S=\|\psi_S\|^2$.
For the unconditioned hard laws define
\begin{equation}
 K_{ST}=\tfrac12(\E_{\widetilde P_+}-\E_{\widetilde P_-})
                         \chi_{S\triangle T}.
 \label{eq:moment}
\end{equation}
The map $|S\rangle\mapsto\psi_S/\sqrt{p_S}$ is an isometry on nonzero
weights. Hence half the difference of the averaged purified output
states has trace norm exactly
\begin{equation}
                  \|\diag(\sqrt p)K\diag(\sqrt p)\|_1 .
 \label{eq:trace}
\end{equation}
Discarding a purifier or measuring cannot increase this norm.
Convexity handles sector mixtures. This reduction includes all repeated
occupations, without identifying photon number with parity size.

The independent permutation signs force adjacent block degrees in
\eqref{eq:moment} to have the same parity; first-block reflection selects
the odd case. Thus every surviving symmetric-difference profile has
odd entries. Define, block by block,
\[
 n_b=|S_b|,\quad m_b=|T_b|,\quad
 a_b=|S_b\triangle T_b|,\quad t_b=|S_b\setminus T_b|.
\]
These are counts of parity-support labels. Since $|S|=|T|$ in one sector,
\begin{equation}
                 m=n+a-2t,\qquad 2|t|=|a|.
 \label{eq:compat}
\end{equation}
Here $4\le|a|\le12$, and independent noise contributes $\beta^{|a|}$.

\section{The coefficient hypothesis and its spectral consequence}\label{sec:ledger}
For an odd profile $a$ and balanced split $t$, let the row labels $A_b$
and column labels $B_b$ be sets of sizes $t_b$ and $a_b-t_b$.
For the exact positive plant $z$ in \eqref{eq:plant}, define the
\emph{complete} marked kernel
\[
 \mathcal O_{a,t}(A,B)=
 \mathbf1_{\{A_b\cap B_b=\varnothing\ \forall b\}}
               \E_z\chi_{A\cup B}(z).
\]
In particular, every distinctness constraint across the split is part
of the matrix. The required analytic hypothesis is
\begin{equation}
 \sup_{p,w}
 \|\diag(\sqrt p)\mathcal O_{a,t}\diag(\sqrt w)\|_1
       \le\widehat c_{a,t},
 \label{eq:coefficient}
\end{equation}
for arbitrary normalized row and column laws $p,w$, with the rational
coefficients used in the stored ledger.

Here is the full lifting argument. Fix $n,m,a,t$ and write $u=a-t$.
For each residual support $R$ of block counts $n-t=m-u$, restrict
$A,B$ to avoid $R$ and embed them as $S=R\cup A$, $T=R\cup B$.
The sum of these embedded matrices is the desired occurrence matrix:
a nonzero entry has the unique residual $R=S\cap T$.
For normalized laws $p_S,w_T$, put
$P_R=\sum_Ap_{R\cup A}$ and $W_R=\sum_Bw_{R\cup B}$.
The contribution of each residual has weighted trace norm at most
$\widehat c_{a,t}\sqrt{P_RW_R}$ by \eqref{eq:coefficient}.
Triangle inequality and Cauchy--Schwarz give
\[
 \widehat c_{a,t}\sum_R\sqrt{P_RW_R}
 \le\widehat c_{a,t}\sqrt{b(n,t)b(m,u)},\qquad
 b(n,t)=\prod_{j=1}^4{n_j\choose t_j}.
\]
Indeed each $S$ appears in exactly $b(n,t)$ residuals, and each $T$
in $b(m,u)$ residuals. This proves the combinatorial factor without
omitting shared labels or distinctness masks.

Index a matrix $A$ by
$\mathcal N=\{n\in\mathbb Z_{\ge0}^4:|n|\le6\}$, which has
${10\choose4}=210$ states. For every compatible profile--split pair
$(a,t)$ with $2|t|=|a|$, add
\begin{equation}
 \beta^{|a|}\widehat c_{a,t}
                 \sqrt{b(n,t)b(m,a-t)}
 \label{eq:edge}
\end{equation}
to $A_{nm}$ and set $B=(A+A^{\mathsf T})/2$.
Restoring the probability of each block-count sector, the preceding
bound gives \eqref{eq:trace}$\le v^{\mathsf T}Bv$, where
$v_n=\sqrt{\Pr(n)}$ and $\|v\|_2=1$.
It is therefore at most $\lambda_{\max}(B)$.
The matrix is nonnegative and block diagonal in total parity-support size.

The ledger reconstructs $B$ with directed decimal arithmetic and
outward square roots. A stored positive vector $w$ provides the bound
\[
 \lambda_{\max}(B)\le\max_i\frac{(\widehat Bw)_i}{w_i}
           <0.258744096385577223.
\]
A floating eigensolve supplies only the candidate; the directed Collatz
ratios supply the spectral certificate for the recorded matrix.
Adding Lemma~\ref{lem:promise} gives \eqref{eq:total} and the Bayes identity
$(1-\TV)/2$ proves the conditional theorem. The total is rounded after
adding the high-precision components; adding their displayed upper
bounds would be slightly looser.

The 888 accepted high-sector rows and inherited lower-degree rows are
mapped to their derivations in \path{code/COEFFICIENTS.md}.
For a nonzero binary64 source value $c$, the assembler uses
$\widehat c=(\lceil10^9c\rceil+1)/10^9$; zero stays zero.
This rounds the supplied value outward, but does not prove
\eqref{eq:coefficient}. Some supporting intermediate bounds are false,
and others lack the stated joint contraction. A new analytic
justification or corrected coefficients are needed before this matrix
certifies a physical lower bound.

The broader fixed-photon-number and adaptive claims also remain open.
The former misses unbalanced parity cuts; the latter additionally needs
a valid contraction retaining the causal measurement structure.
See \path{../AUDIT.md} for exact counterexamples to the supplied
inferences and \path{../REVIEW.md} for the independent reviews.
Run \path{make check} to reproduce the arithmetic and counterexamples,
and \path{make build} to build both proof notes.

\bibliographystyle{plain}
\bibliography{references}
\end{document}
